\documentclass[uplatex,dvipdfmx,a4paper,11pt]{jsarticle}
\input{../../../templates/kaisetsu.tex}
\kaisetsuprint{連続体力学　第3回「応力テンソルと歪みテンソル」}{2026-10-28}
%% kaisetsu-moto	problems\連続体力学\kai03\連続体力学-3.tex	kiroku\kaisetsu\連続体力学\連続体力学-3_2026-10-28_komento.tex
\begin{document}
\matome
\begin{saitenichiran}
  \hyougyou{1}{L1・計算}{◎◎◎◎}{S}{全問 ◎}
  \hyougyou{\hyperlink{bt-m-019-rigid}{2}}{L1・導出}{○◎○}{B}{(1)記述 (3)記述}
  \hyougyou{\hyperlink{bt-m-020-sym}{3}}{L1・計算}{◎◎△}{B}{(3)計算}
  \hyougyou{\hyperlink{bt-m-021-types}{4}}{L1・説明}{◎◎○△}{B}{(3)記述 (4)方針}
  \hyougyou{\hyperlink{bt-m-022-decomp}{5}}{L2・導出}{◎△◎◎}{A}{(2)記述}
  \hyougyou{6}{L2・計算}{◎◎◎◎}{S}{全問 ◎}
  \hyougyou{\hyperlink{bt-m-024-shear}{7}}{L2・計算}{◎◎△}{B}{(3)方針}
\end{saitenichiran}
\begin{daiji}
  \item 回転角は $\vect{\Omega}$ の大きさの半分である（問題4 (4)・問題7 (3)）。問題1 (4) ではできていたので、$\vect{\Omega}$ を求めたら「角度はその半分」と書く手順を毎回入れる。
  \item 接線成分は $\vect{p}-(\vect{p}\cdot\vect{n})\vect{n}$ を最後まで引き算し、$\abs{\vect{p}}^{2}=p_{n}^{2}+\abs{\vect{p}_{t}}^{2}$ で検算する（問題3 (3)）。
  \item 「ただ一通りに分けられる」は、別の分け方があると仮定して差が 0 になることで示す（問題5 (2)）。
\end{daiji}
\begin{maemo}
  \item 回転の向き・角度の取り違え：第1回 問題6 (4)（2026-10-21）で等密度面の回転の向きを逆にした。回転は、具体的な1点の変位を計算して向きと角度を確かめる。
\end{maemo}

\youtenseiri
\begin{noteshikilist}
  \item 応力テンソルの対称性 % 出典: p.10
    \[ P_{ij}=P_{ji} \notetag{3} \]
  \item 変形勾配テンソルの成分 % 出典: p.19
    \[ D_{ij}=\pdv{r_i(\vect{x})}{x_j} \notetag{13} \]
  \item 歪みテンソル（対称部分） % 出典: p.20
    \[ E_{ij}=\frac{1}{2}\left(\pdv{r_i}{x_j}+\pdv{r_j}{x_i}\right) \notetag{16} \]
  \item 回転テンソル（反対称部分） % 出典: p.20
    \[ F_{ij}=\frac{1}{2}\left(\pdv{r_i}{x_j}-\pdv{r_j}{x_i}\right) \notetag{17} \]
  \item 回転ベクトルの成分 % 出典: p.33
    \[ 2F_{32}=(\vect{\nabla}\times\vect{r})_1=\Omega_1 \notetag{33} \]
    \[ 2F_{13}=(\vect{\nabla}\times\vect{r})_2=\Omega_2 \notetag{34} \]
    \[ 2F_{21}=(\vect{\nabla}\times\vect{r})_3=\Omega_3 \notetag{35} \]
  \item 回転テンソルの回転ベクトルによる表示 % 出典: p.34
    \[ \tens{F}=\begin{pmatrix} 0 & -\frac{1}{2}\Omega_3 & \frac{1}{2}\Omega_2\\ \frac{1}{2}\Omega_3 & 0 & -\frac{1}{2}\Omega_1\\ -\frac{1}{2}\Omega_2 & \frac{1}{2}\Omega_1 & 0 \end{pmatrix} \notetag{37} \]
\end{noteshikilist}

\mondaihen
\level{1}
\mondaibangou{2}
\begin{mondai}{019-rigid}{導出}{平行移動は変形を生まず，微小な剛体回転は回転テンソルだけを持つこと（回転角は $\Omega$ の半分）}
定ベクトル $\vect{c}$ と定数 $\theta$（$\abs{\theta}\ll1$）を用いて，変位場
  \[ \vect{r}(\vect{x})=\vect{c}+\theta\,\vect{e}_3\times\vect{x} \]
  を考える。
  \begin{shoumon}
    \item $\vect{r}$ の成分を書き，近接2点の変形 $\delta\vect{r}=\vect{r}(\vect{x}+\delta\vect{x})-\vect{r}(\vect{x})$ を $\delta\vect{x}$ の成分で表せ。$\vect{c}$ が $\delta\vect{r}$ に現れないことを示し，その物理的な意味を述べよ。
    \item $\tens{D}$，$\tens{E}$，$\tens{F}$ と回転ベクトル $\vect{\Omega}$ を求め，$x_3$ 軸まわりの回転角と $\Omega_3$ の関係を示せ。
    \item 有限の角度 $\theta$ の $x_3$ 軸まわりの回転では，点 $\vect{x}$ は $(x_1\cos\theta-x_2\sin\theta,\ x_1\sin\theta+x_2\cos\theta,\ x_3)$ へ移る。この変位場について $\tens{E}$ を求め，その0でない成分が $\theta$ について何次の微小量かを答えよ。そこから，(2) の結論が成り立つ条件を述べよ。
  \end{shoumon}

\begin{kotaeran}
  \machigai{1}{○}{記述} $\delta\vect{r}=(-\theta\delta x_{2},\theta\delta x_{1},0)$ で $\vect{c}$ が消えることは正しい。問われている $\vect{c}$ の物理的な意味が「$\delta\vect{r}$ は変位だから」で途切れている。$\vect{c}$ は全部の点を同じだけ動かす平行移動で、2点の相対位置 $\delta\vect{r}$ を変えないので変形を生まない、まで書く。
  \machigai{3}{○}{記述} 展開を $\cos\theta-1=1-\theta^{2}/2+\dotsb$ と書いている。これは $\cos\theta$ の展開で、正しくは $\cos\theta-1=-\theta^{2}/2+\dotsb$（$\theta$ の2次）。
\end{kotaeran}
\begin{rikai}
  \item 変形を表すのは、近くの2点の相対位置 $\delta\vect{r}$ の変化である。全体の平行移動 $\vect{c}$ は $\delta\vect{r}$ に入らないので、変形とは無関係。
  \item 回転による長さの変化は $\theta$ の2次なので、$\theta\ll1$ では歪み $\tens{E}$ は 0 として扱える。
\end{rikai}
\end{mondai}

\mondaibangou{3}
\begin{mondai}{020-sym}{計算}{応力テンソルの対称性と，応力ベクトルの法線成分・接線成分}
ある点の応力テンソルの成分が，定数 $p_0,\tau,\sigma$ を用いて
  \[ \tens{P}=\begin{pmatrix} -p_0 & \tau & P_{13}\\ P_{21} & -p_0 & 0\\ \sigma & P_{32} & -p_0 \end{pmatrix} \]
  と分かっている（$P_{13},P_{21},P_{32}$ は未知）。$p_0$ は圧力の大きさを表す定数で，応力ベクトルの成分 $p_i$ とは別の量である。応力ベクトルは $\vect{p}(\vect{x},\vect{n})=\tens{P}\cdot\vect{n}$ で与えられる。
  \begin{shoumon}
    \item 応力テンソルの対称性を用いて $P_{13},P_{21},P_{32}$ を求めよ。また，一般の応力テンソルの独立な成分は何個か。
    \item 単位法線ベクトル $\vect{n}=(\vect{e}_1+\vect{e}_2)/\sqrt{2}$ の面に働く応力ベクトル $\vect{p}$ を求めよ。
    \item $\vect{p}$ の法線方向の成分 $\vect{p}\cdot\vect{n}$ と，面に平行な成分の大きさを求めよ。
  \end{shoumon}

\begin{kotaeran}
  \machigai{3}{△}{計算} 法線成分 $\tau-p_{0}$ と、接線成分を $\vect{p}-(\vect{p}\cdot\vect{n})\vect{n}$ で出す方針は正しい。引いた結果を $\bigl((\tau-p_{0})/\sqrt{2},(\tau-p_{0})/\sqrt{2},\sigma\bigr)$ と書いていて、これは $\vect{p}$ そのまま（第3成分の $1/\sqrt{2}$ も落ちている）。$(\vect{p}\cdot\vect{n})\vect{n}=\bigl((\tau-p_{0})/\sqrt{2},(\tau-p_{0})/\sqrt{2},0\bigr)$ を引くと第1・第2成分は消えて $(0,0,\sigma/\sqrt{2})$、大きさは $\abs{\sigma}/\sqrt{2}$。
\end{kotaeran}
\begin{rikai}
  \item 応力ベクトルの分け方：法線成分 $p_{n}=\vect{p}\cdot\vect{n}$、接線ベクトル $\vect{p}_{t}=\vect{p}-p_{n}\vect{n}$。$\vect{p}_{t}$ は $\vect{n}$ に垂直なので、$\vect{p}_{t}\cdot\vect{n}=0$ になっているかで引き算を確かめられる。
  \item 検算：$\abs{\vect{p}}^{2}=(\tau-p_{0})^{2}+\sigma^{2}/2$ が $p_{n}^{2}+\abs{\vect{p}_{t}}^{2}$ に等しいはず。答え $\sqrt{(\tau-p_{0})^{2}+\sigma^{2}}$ だと、これより大きくなってしまう。
\end{rikai}
\end{mondai}

\mondaibangou{4}
\begin{mondai}{021-types}{説明}{変形勾配テンソルの形から変形の種類を読み取り，図示する}
定数 $\varepsilon,\gamma,\theta$ はいずれも正で $1$ より十分小さい。次の4つの変形勾配テンソルを考える（書いていない成分は0）。
  \[ \text{(a)}\ D_{22}=-\varepsilon,\qquad \text{(b)}\ D_{13}=D_{31}=\frac{\gamma}{2},\qquad \text{(c)}\ D_{32}=\theta,\ D_{23}=-\theta,\qquad \text{(d)}\ D_{11}=D_{22}=D_{33}=\varepsilon \]
  \begin{shoumon}
    \item (a)〜(d) のそれぞれについて，$\tens{E}$ と $\tens{F}$ が0かどうかを答え，どのような変形か（伸縮・剪断・回転など）を述べよ。
    \item (a) について，$x_1$--$x_2$ 面内で一辺 $\ell$ の正方形（2辺が $x_1$ 軸・$x_2$ 軸に沿い，1つの頂点が原点）が変形後にどうなるかを図示せよ。原点は動かないものとする。
    \item (b) について，$x_1$--$x_3$ 面内で同様の正方形が変形後にどうなるかを図示せよ。
    \item (c) について，回転ベクトル $\vect{\Omega}$ を求め，どの軸のまわりに，どちら向きに，どれだけ回転するかを答えよ。
  \end{shoumon}

\begin{kotaeran}
  \machigai{3}{○}{記述} ひし形の形は正しい。図の頂点を $(\gamma/2,\ell)$、$(\ell,\gamma/2)$ と書いたが、正しくは $(\gamma\ell/2,\ell)$、$(\ell,\gamma\ell/2)$。座標は長さなので、無次元の $\gamma/2$ だけでは次元が合わない。
  \machigai{4}{△}{方針} $\vect{\Omega}=(2\theta,0,0)$ と $x_{1}$ 軸まわりの反時計回りは正しいが、回転角を $2\theta$ とした。回転角は $\Omega_{1}/2=\theta$。
\end{kotaeran}
\begin{rikai}
  \item 回転を表す部分は $\tens{F}\cdot\delta\vect{x}=\frac{1}{2}\vect{\Omega}\times\delta\vect{x}$ で、$\frac{1}{2}$ が付いている。角度 $\theta$ の微小回転は $\delta\vect{r}=\vect{\theta}\times\delta\vect{x}$ なので、$\vect{\theta}=\vect{\Omega}/2$、つまり回転角は $\abs{\vect{\Omega}}$ の半分になる。
  \item 確かめ方：$\tens{D}$ から $\delta\vect{r}=(0,-\theta\delta x_{3},\theta\delta x_{2})$ を作ると、角度 $\theta$ の回転そのものになっている。
  \item 変位・座標は「歪み（無次元）$\times$ 長さ」。図に座標を書いたら次元を確かめる。
\end{rikai}
\end{mondai}

\level{2}
\mondaibangou{5}
\begin{mondai}{022-decomp}{導出}{変形勾配テンソルの導出，対称・反対称分解の一意性，回転テンソルの作用が回転ベクトルとの外積になること}
変位場 $\vect{r}(\vect{x})$ は滑らかで，変形は微小であるとする。
  \begin{shoumon}
    \item 近接2点の変形 $\delta\vect{r}=\vect{r}(\vect{x}+\delta\vect{x})-\vect{r}(\vect{x})$ を $\delta\vect{x}$ について1次まで展開し，$\delta\vect{r}\simeq\tens{D}\cdot\delta\vect{x}$ となること，$D_{ij}$ が1点 $\vect{x}$ だけで決まることを示せ。
    \item 任意の $\tens{D}$ は，対称テンソル $\tens{E}$ と反対称テンソル $\tens{F}$ の和 $\tens{D}=\tens{E}+\tens{F}$ にただ一通りに分けられることを示し，$\tens{E}$，$\tens{F}$ を $\tens{D}$ とその転置 $\tens{D}^{\mathrm{T}}$（成分 $(\tens{D}^{\mathrm{T}})_{ij}=D_{ji}$）で表せ。
    \item 回転ベクトル $\vect{\Omega}=\vect{\nabla}\times\vect{r}$ を用いて，任意の $\delta\vect{x}$ に対して
    \[ \tens{F}\cdot\delta\vect{x}=\frac{1}{2}\vect{\Omega}\times\delta\vect{x} \]
    が成り立つことを成分で示せ。
    \item $\tens{E}$ と $\tens{F}$ の独立な成分はそれぞれ何個か。(3) と合わせて，$\tens{F}$ の独立な成分が何で表されるかを述べよ。
  \end{shoumon}

\begin{kotaeran}
  \machigai{2}{△}{記述} $\tens{E}=\frac{1}{2}(\tens{D}+\tens{D}^{\mathrm{T}})$、$\tens{F}=\frac{1}{2}(\tens{D}-\tens{D}^{\mathrm{T}})$ は正しいが、問われている「ただ一通りに分けられること」の証明が無い。
\end{kotaeran}
\begin{rikai}
  \item 一意性の示し方：別の分け方 $\tens{D}=\tens{E}'+\tens{F}'$（$\tens{E}'$ 対称・$\tens{F}'$ 反対称）があると仮定すると、$\tens{E}-\tens{E}'=\tens{F}'-\tens{F}$ になる。左辺は対称、右辺は反対称なので、この行列 $\tens{A}$ は $\tens{A}=\tens{A}^{\mathrm{T}}=-\tens{A}$ を満たし $\tens{A}=0$。よって $\tens{E}'=\tens{E}$、$\tens{F}'=\tens{F}$。
  \item 分け方を1つ見つけること（存在）と、それしか無いこと（一意性）は別の主張である。
\end{rikai}
\end{mondai}

\mondaibangou{7}
\begin{mondai}{024-shear}{計算}{単純剪断が純粋剪断と回転の和であること（逆向きの組み立て）と図示}
$x_2$--$x_3$ 面内の純粋剪断を表す歪みテンソル
  \[ E_{23}=E_{32}=\frac{\gamma}{2}\quad(\text{他の成分は}0,\ 0<\gamma\ll1) \]
  と，$x_1$ 軸方向の回転ベクトル $\vect{\Omega}=\Omega_1\vect{e}_1$ に対応する回転テンソル $\tens{F}$ を足し合わせて，変形勾配テンソル $\tens{D}=\tens{E}+\tens{F}$ を作る。
  \begin{shoumon}
    \item $\tens{F}$ の成分を $\Omega_1$ で表し，$\tens{D}$ を求めよ。
    \item $\tens{D}$ の0でない成分がちょうど1つになるような $\Omega_1$ をすべて求めよ。それぞれについて $\tens{D}$ と変形 $\delta\vect{r}$ の成分を書け。
    \item (2) のうち $\Omega_1>0$ の場合について，$x_2$--$x_3$ 面内で一辺 $\ell$ の正方形（2辺が $x_2$ 軸・$x_3$ 軸に沿い，1つの頂点が原点。原点は動かない）が変形後にどうなるかを図示せよ。また，この変形に含まれる回転の向きと角度を答えよ。
  \end{shoumon}

\begin{kotaeran}
  \machigai{3}{△}{方針} 図（$x_{3}$ 軸に沿う辺はそのままで、$x_{2}$ 軸に沿う辺が $x_{3}$ 側へ傾く平行四辺形）と回転の向きは正しいが、角度を $\Omega_{1}$ とした。正しくは $\Omega_{1}/2=\gamma/2$。途中の図の頂点も $(\ell,\gamma/2)$ と $\ell$ が抜けている（問題4 (3) と同じ）。
\end{kotaeran}
\begin{rikai}
  \item 問題4 (4) と同じく、回転角は $\abs{\vect{\Omega}}$ の半分。単純ずりは「角度 $\gamma/2$ の回転」と「純粋な歪み」に分かれる（$\gamma$ 全部が回転ではない）。
\end{rikai}
\end{mondai}

\kaitouhen
\level{1}
\kaisetsubangou{2}
\begin{kaisetsu}{019-rigid}
\houshin 変形の定義から $\delta\vect{r}$ を直接計算し，\noteshiki{13}・\noteshiki{16}・\noteshiki{17}・\noteshiki{35} で各テンソルと回転ベクトルを求める。(3) は有限回転の変位場で $\tens{E}$ を計算する。
\keisan \shoumonkai{1} $\vect{e}_3\times\vect{x}=(-x_2,x_1,0)$ なので
\[ \vect{r}(\vect{x})=(c_1-\theta x_2,\ c_2+\theta x_1,\ c_3). \]
差をとると，定数 $c_1,c_2,c_3$ は打ち消し合い
\[ \delta\vect{r}=\vect{r}(\vect{x}+\delta\vect{x})-\vect{r}(\vect{x})=(-\theta\,\delta x_2,\ \theta\,\delta x_1,\ 0) \]
となる（変位場が $\vect{x}$ の1次式なので，この式は近似なしに成り立つ）。$\vect{c}$ は全点を同じだけ動かす平行移動で，2点の変位の差を変えないので変形を生まない。

\shoumonkai{2} 0 でない成分は $D_{12}=\pdv{r_1}{x_2}=-\theta$ と $D_{21}=\pdv{r_2}{x_1}=\theta$。$\tens{D}$ は反対称なので
\[ \tens{E}=0,\qquad \tens{F}=\tens{D}=\begin{pmatrix} 0 & -\theta & 0\\ \theta & 0 & 0\\ 0 & 0 & 0 \end{pmatrix}. \]
\noteshiki{35} より $\Omega_3=2F_{21}=2\theta$，$\Omega_1=\Omega_2=0$。よって $\vect{\Omega}=2\theta\,\vect{e}_3$ で，$x_3$ 軸まわりの回転角 $\theta$ は $\Omega_3/2$ に等しい。

\shoumonkai{3} 変位は移った先の位置から元の位置を引いたものなので
\[ r_1=x_1(\cos\theta-1)-x_2\sin\theta,\qquad r_2=x_1\sin\theta+x_2(\cos\theta-1),\qquad r_3=0. \]
偏微分すると $D_{11}=D_{22}=\cos\theta-1$，$D_{12}=-\sin\theta$，$D_{21}=\sin\theta$，他は0。対称部分は
\[ E_{11}=E_{22}=\cos\theta-1=-\frac{\theta^{2}}{2}+O(\theta^{4}),\qquad E_{12}=\frac{-\sin\theta+\sin\theta}{2}=0 \]
で，0 でない成分は $\theta$ の2次の微小量である。したがって「剛体回転では $\tens{E}=0$」は $\theta$ の1次まで（微小変形の範囲）で成り立つ結論であり，$\theta^{2}$ が無視できない大きな回転には使えない。
\kotae{\shoumonkai{1} $\delta\vect{r}=(-\theta\,\delta x_2,\ \theta\,\delta x_1,\ 0)$。平行移動 $\vect{c}$ は変形を生まない \shoumonkai{2} $D_{12}=-\theta$，$D_{21}=\theta$；$\tens{E}=0$，$\tens{F}=\tens{D}$；$\vect{\Omega}=2\theta\,\vect{e}_3$，回転角 $\theta=\Omega_3/2$ \shoumonkai{3} $E_{11}=E_{22}=\cos\theta-1\simeq-\theta^{2}/2$（2次），他は0。$\abs{\theta}\ll1$ で2次を無視できるとき}
\kenzan $\theta\to0$ で $\tens{D}\to0$（変形なし）。(3) で $\theta$ の1次までとると $D_{12}=-\theta$，$D_{21}=\theta$，$D_{11}=D_{22}=0$ となり (2) に一致する。$\delta\vect{r}\cdot\delta\vect{x}=-\theta\,\delta x_1\delta x_2+\theta\,\delta x_2\delta x_1=0$ で，点は回転の円周方向に動いている。
\tsumazuki 有限回転の $\tens{E}$ が0にならないのは，変形勾配から歪みを $\tens{E}=(\tens{D}+\tens{D}^{\mathrm{T}})/2$ と定めたのが微小変形の近似だからである。
\end{kaisetsu}

\kaisetsubangou{3}
\begin{kaisetsu}{020-sym}
\houshin 応力テンソルの対称性 \noteshiki{3} で未知成分を決め，$\vect{p}=\tens{P}\cdot\vect{n}$ を成分で計算する。法線成分は内積，接線成分は残りのベクトルの大きさ。
\keisan \shoumonkai{1} $P_{13}=P_{31}=\sigma$，$P_{21}=P_{12}=\tau$，$P_{32}=P_{23}=0$。
\[ \tens{P}=\begin{pmatrix} -p_0 & \tau & \sigma\\ \tau & -p_0 & 0\\ \sigma & 0 & -p_0 \end{pmatrix}. \]
対称性により，9個の成分のうち独立なのは対角3個と非対角3個の6個である。

\shoumonkai{2} $n_1=n_2=1/\sqrt{2}$，$n_3=0$ なので
\[ p_i=\sum_{j=1}^{3}P_{ij}n_j=\frac{1}{\sqrt{2}}\left(P_{i1}+P_{i2}\right). \]
各成分は
\[ p_1=\frac{-p_0+\tau}{\sqrt{2}},\qquad p_2=\frac{\tau-p_0}{\sqrt{2}},\qquad p_3=\frac{\sigma+0}{\sqrt{2}}. \]
\shoumonkai{3} 法線成分は
\[ \vect{p}\cdot\vect{n}=\frac{1}{2}\bigl[(\tau-p_0)+(\tau-p_0)\bigr]=\tau-p_0. \]
面に平行な成分は
\[ \vect{p}-(\vect{p}\cdot\vect{n})\vect{n}=\frac{1}{\sqrt{2}}(\tau-p_0,\ \tau-p_0,\ \sigma)-\frac{\tau-p_0}{\sqrt{2}}(1,\ 1,\ 0)=\left(0,\ 0,\ \frac{\sigma}{\sqrt{2}}\right) \]
で，大きさは $\abs{\sigma}/\sqrt{2}$。
\kotae{\shoumonkai{1} $P_{13}=\sigma$，$P_{21}=\tau$，$P_{32}=0$。独立な成分は6個 \shoumonkai{2} $\vect{p}=\frac{1}{\sqrt{2}}(\tau-p_0,\ \tau-p_0,\ \sigma)$ \shoumonkai{3} 法線成分 $\tau-p_0$，面に平行な成分の大きさ $\abs{\sigma}/\sqrt{2}$}
\kenzan $\abs{\vect{p}}^{2}=(\tau-p_0)^{2}+\sigma^{2}/2$ は，法線成分の2乗 $(\tau-p_0)^{2}$ と接線成分の2乗 $\sigma^{2}/2$ の和に等しい。$\tau=\sigma=0$（静止流体，$\tens{P}=-p_0\tens{I}$）なら $\vect{p}=-p_0\vect{n}$ で，圧力は面を押す向きに働く。
\tsumazuki $\vect{p}\cdot\vect{n}=\tau-p_0$ は，圧力 $p_0$ だけでなく剪断成分 $\tau$ も含む。面の向きを変えると，非対角成分が法線成分に入ってくる。
\end{kaisetsu}

\kaisetsubangou{4}
\begin{kaisetsu}{021-types}
\houshin $\tens{D}$ の対称・反対称を見て $\tens{E}$，$\tens{F}$ を判定し，$\delta\vect{r}=\tens{D}\cdot\delta\vect{x}$ で正方形の頂点の移動を計算する。
\keisan \shoumonkai{1} (a) 対角成分だけなので対称：$\tens{F}=0$，$\tens{E}=\tens{D}\neq0$。$\delta r_2=-\varepsilon\,\delta x_2$ で，$x_2$ 方向だけの一軸の縮み（圧縮）。(b) $D_{13}=D_{31}$ で対称：$\tens{F}=0$，$\tens{E}\neq0$。$x_1$--$x_3$ 面内の純粋剪断。(c) $D_{23}=-D_{32}$ で反対称：$\tens{E}=0$，$\tens{F}=\tens{D}\neq0$。$x_1$ 軸まわりの回転。(d) 対角で対称：$\tens{F}=0$。3軸方向に同じ割合 $\varepsilon$ で伸びる，等方的な膨張（多軸伸長で3軸が等しい場合）。

\shoumonkai{2} 頂点 $(\ell,0)$ は動かず，$(0,\ell)$ と $(\ell,\ell)$ は $\delta r_2=-\varepsilon\ell$ だけ下がる。正方形は，横 $\ell$・縦 $(1-\varepsilon)\ell$ の長方形になる（下図左）。

\shoumonkai{3} $(x_1,x_3)$ 座標で，$\delta r_1=\frac{\gamma}{2}\delta x_3$，$\delta r_3=\frac{\gamma}{2}\delta x_1$。頂点 $(\ell,0)$ は $(\ell,\frac{\gamma}{2}\ell)$ へ，$(0,\ell)$ は $(\frac{\gamma}{2}\ell,\ell)$ へ，$(\ell,\ell)$ は $\bigl((1+\gamma/2)\ell,\ (1+\gamma/2)\ell\bigr)$ へ移る。$x_1$ 軸に沿う辺は $x_3$ 側へ，$x_3$ 軸に沿う辺は $x_1$ 側へ，同じ角度 $\gamma/2$ ずつ傾き，正方形は対角線について対称なひし形になる（下図右）。
\zu[(a) $x_2$ 方向の圧縮（左，$x_1$--$x_2$ 面）と (b) 純粋剪断（右，$x_1$--$x_3$ 面）。破線が変形前，実線が変形後（変形は誇張して描いた）]{
 \begin{tikzpicture}
  \draw[jiku] (0,0) -- (2.4,0) node[right] {$x_1$};
  \draw[jiku] (0,0) -- (0,2.3) node[above] {$x_2$};
  \draw[hojo] (0,0) rectangle (1.6,1.6);
  \draw[thick] (0,0) rectangle (1.6,1.2);
  \node[below] at (0.8,0) {$\ell$};
  \node[right] at (1.6,0.6) {$(1-\varepsilon)\ell$};
  \begin{scope}[xshift=5cm]
    \draw[jiku] (0,0) -- (2.6,0) node[right] {$x_1$};
    \draw[jiku] (0,0) -- (0,2.6) node[above] {$x_3$};
    \draw[hojo] (0,0) rectangle (1.6,1.6);
    \draw[thick] (0,0) -- (1.6,0.4) -- (2.0,2.0) -- (0.4,1.6) -- cycle;
    \node[right] at (1.6,0.4) {$(\ell,\frac{\gamma}{2}\ell)$};
    \node[above] at (0.4,1.6) {$(\frac{\gamma}{2}\ell,\ell)$};
  \end{scope}
 \end{tikzpicture}
}
\shoumonkai{4} $F_{32}=(D_{32}-D_{23})/2=\theta$，$F_{13}=F_{21}=0$ なので，\noteshiki{33}--\noteshiki{35} より $\vect{\Omega}=(2\theta,0,0)$。$\delta\vect{r}=(0,-\theta\,\delta x_3,\theta\,\delta x_2)$ で，$x_2$ 軸上の点は $+x_3$ 向き，$x_3$ 軸上の点は $-x_2$ 向きに動く。$x_1$ 軸まわりに，$x_2$ 軸が $x_3$ 軸に向かう向き（$+x_1$ から見て反時計回り）に角度 $\Omega_1/2=\theta$ の回転である。
\kotae{\shoumonkai{1} (a) $\tens{F}=0$：$x_2$ 方向の一軸圧縮 (b) $\tens{F}=0$：$x_1$--$x_3$ 面の純粋剪断 (c) $\tens{E}=0$：$x_1$ 軸まわりの回転 (d) $\tens{F}=0$：等方的な膨張 \shoumonkai{2} 横 $\ell$・縦 $(1-\varepsilon)\ell$ の長方形 \shoumonkai{3} 2辺が互いに $\gamma/2$ ずつ傾いたひし形 \shoumonkai{4} $\vect{\Omega}=(2\theta,0,0)$。$x_1$ 軸まわりに $x_2\to x_3$ の向きへ角度 $\theta$}
\kenzan (c) で $\delta\vect{r}\cdot\delta\vect{x}=-\theta\,\delta x_2\delta x_3+\theta\,\delta x_3\delta x_2=0$ となり，長さを変えない回転と整合する。(b) の図は対角線 $x_1=x_3$ について対称で，$\tens{E}$ が対称であることと合う。
\tsumazuki (b) の純粋剪断では2辺が互いに逆向きに傾く。一方の辺だけが傾く単純剪断（$D_{13}$ だけが0でない場合など）とは別の変形である。
\end{kaisetsu}

\level{2}
\kaisetsubangou{5}
\begin{kaisetsu}{022-decomp}
\houshin (1) は Taylor 展開，(2) は「対称かつ反対称なテンソルは0」を使う。(3) は \noteshiki{37} の $\tens{F}$ を $\delta\vect{x}$ に掛けて外積と成分ごとに比べる。
\keisan \shoumonkai{1} 各成分 $r_i$ を $\vect{x}$ のまわりで展開すると
\[ r_i(\vect{x}+\delta\vect{x})=r_i(\vect{x})+\sum_{j=1}^{3}\pdv{r_i(\vect{x})}{x_j}\delta x_j+O\bigl((\delta x)^{2}\bigr). \]
したがって
\[ \delta r_i=\sum_{j=1}^{3}D_{ij}\,\delta x_j+O\bigl((\delta x)^{2}\bigr),\qquad D_{ij}=\pdv{r_i(\vect{x})}{x_j}. \]
2次以上を無視すると $\delta\vect{r}\simeq\tens{D}\cdot\delta\vect{x}$ である。$D_{ij}$ は点 $\vect{x}$ での変位場の微分であり，相対位置 $\delta\vect{x}$ に依らない。よって $\tens{D}$ は1点 $\vect{x}$ の関数（場）になる。

\shoumonkai{2} 存在：$\tens{E}=(\tens{D}+\tens{D}^{\mathrm{T}})/2$，$\tens{F}=(\tens{D}-\tens{D}^{\mathrm{T}})/2$ とおくと，$E_{ji}=E_{ij}$，$F_{ji}=-F_{ij}$ で，和は $\tens{D}$ になる。

一意性：別の分け方 $\tens{D}=\tens{E}'+\tens{F}'$（$\tens{E}'$ 対称，$\tens{F}'$ 反対称）があるとすると
\[ \tens{E}-\tens{E}'=\tens{F}'-\tens{F}. \]
左辺は対称，右辺は反対称なので，この量 $\tens{G}$ は $G_{ij}=G_{ji}$ と $G_{ij}=-G_{ji}$ を同時に満たし，$G_{ij}=0$。よって $\tens{E}'=\tens{E}$，$\tens{F}'=\tens{F}$ で，分け方はただ一通りである。

\shoumonkai{3} \noteshiki{37} より，第1成分は
\[ (\tens{F}\cdot\delta\vect{x})_1=F_{12}\,\delta x_2+F_{13}\,\delta x_3=-\frac{1}{2}\Omega_3\,\delta x_2+\frac{1}{2}\Omega_2\,\delta x_3. \]
外積の第1成分は
\[ \left(\frac{1}{2}\vect{\Omega}\times\delta\vect{x}\right)_1=\frac{1}{2}\left(\Omega_2\,\delta x_3-\Omega_3\,\delta x_2\right) \]
で，一致する。第2・第3成分も同様に計算すると
\begin{align*}
(\tens{F}\cdot\delta\vect{x})_2&=F_{21}\,\delta x_1+F_{23}\,\delta x_3=\frac{1}{2}\left(\Omega_3\,\delta x_1-\Omega_1\,\delta x_3\right)=\left(\frac{1}{2}\vect{\Omega}\times\delta\vect{x}\right)_2,\\
(\tens{F}\cdot\delta\vect{x})_3&=F_{31}\,\delta x_1+F_{32}\,\delta x_2=\frac{1}{2}\left(\Omega_1\,\delta x_2-\Omega_2\,\delta x_1\right)=\left(\frac{1}{2}\vect{\Omega}\times\delta\vect{x}\right)_3.
\end{align*}
\shoumonkai{4} 対称な $\tens{E}$ は対角3個と非対角3個で6個，反対称な $\tens{F}$ は対角が0で $F_{21},F_{13},F_{32}$ の3個（合計9個で $\tens{D}$ の成分数と一致）。(3) より，$\tens{F}$ の3個の独立な成分は回転ベクトル $\vect{\Omega}$ の3成分で表され，$\tens{F}$ による変形は角速度ベクトルのような外積 $\frac{1}{2}\vect{\Omega}\times\delta\vect{x}$，すなわち $\vect{\Omega}$ を軸とする角度 $\abs{\vect{\Omega}}/2$ の微小回転である。
\kotae{\shoumonkai{1} $\delta\vect{r}\simeq\tens{D}\cdot\delta\vect{x}$，$D_{ij}=\pdv{r_i(\vect{x})}{x_j}$（$\delta\vect{x}$ に依らない） \shoumonkai{2} $\tens{E}=(\tens{D}+\tens{D}^{\mathrm{T}})/2$，$\tens{F}=(\tens{D}-\tens{D}^{\mathrm{T}})/2$（一意） \shoumonkai{3} 成分ごとに一致（本文） \shoumonkai{4} $\tens{E}$ は6個，$\tens{F}$ は3個で，$\vect{\Omega}$ の3成分で表される}
\kenzan (3) を $x_3$ 軸まわりの回転（$\Omega_3$ だけが0でない場合）に当てはめると
\[ \frac{1}{2}\vect{\Omega}\times\delta\vect{x}=\frac{1}{2}\Omega_3(-\delta x_2,\ \delta x_1,\ 0) \]
で，要点整理の回転の変形 $\delta r_1=-\frac{1}{2}\Omega_3\delta x_2$，$\delta r_2=\frac{1}{2}\Omega_3\delta x_1$ に一致する。次元：$\tens{D}$，$\vect{\Omega}$ は無次元で，両辺とも長さの次元。
\tsumazuki 展開の係数 $\pdv{r_i}{x_j}$ を $\vect{x}+\delta\vect{x}$ で評価しても違いは $\delta\vect{x}$ の1次で，$\delta\vect{r}$ への寄与は2次になるので無視できる。
\end{kaisetsu}

\kaisetsubangou{7}
\begin{kaisetsu}{024-shear}
\houshin \noteshiki{37} で $\tens{F}$ を $\Omega_1$ で書き，$\tens{E}$ に足す。0 でない成分が1つになる条件を課し，$\delta\vect{r}=\tens{D}\cdot\delta\vect{x}$ で図示する。
\keisan \shoumonkai{1} $\Omega_2=\Omega_3=0$ とおくと，\noteshiki{37} より $F_{32}=\Omega_1/2$，$F_{23}=-\Omega_1/2$，他は0。よって
\[ \tens{D}=\begin{pmatrix} 0 & 0 & 0\\ 0 & 0 & \frac{\gamma-\Omega_1}{2}\\ 0 & \frac{\gamma+\Omega_1}{2} & 0 \end{pmatrix}. \]
\shoumonkai{2} $\gamma>0$ なので，$D_{23}$ と $D_{32}$ が同時に0になることはない。どちらか一方を0にする。

$\Omega_1=\gamma$ のとき：$D_{23}=0$，$D_{32}=\gamma$。$\delta\vect{r}=(0,\ 0,\ \gamma\,\delta x_2)$。

$\Omega_1=-\gamma$ のとき：$D_{32}=0$，$D_{23}=\gamma$。$\delta\vect{r}=(0,\ \gamma\,\delta x_3,\ 0)$。

どちらも，一方向の変位がそれと垂直な座標に比例する単純剪断である。

\shoumonkai{3} $\Omega_1=\gamma$ の場合，$(x_2,x_3)$ 座標で頂点 $(\ell,0)$ は $(\ell,\gamma\ell)$ へ，$(\ell,\ell)$ は $(\ell,(1+\gamma)\ell)$ へ移り，$x_3$ 軸上の頂点 $(0,\ell)$ は動かない。$x_3$ 軸に沿う辺はそのままで，$x_2$ 軸に沿う辺が $x_3$ 側へ角度 $\gamma$ だけ傾き，正方形は平行四辺形になる（下図）。含まれる回転は $x_1$ 軸まわりに，$x_2$ 軸が $x_3$ 軸に向かう向き（$+x_1$ から見て反時計回り）の角度 $\Omega_1/2=\gamma/2$ である。これは，2辺がそれぞれ $\gamma/2$ ずつ互いに近づく純粋剪断に，全体を $\gamma/2$ 回す回転を加えると，$x_3$ 軸に沿う辺の傾きは打ち消され，$x_2$ 軸に沿う辺の傾きは $\gamma$ になることに対応する。
\zu[$\Omega_1=\gamma$ の単純剪断（$x_2$--$x_3$ 面）。破線が変形前，実線が変形後（変形は誇張して描いた）]{
 \begin{tikzpicture}
  \draw[jiku] (0,0) -- (2.4,0) node[right] {$x_2$};
  \draw[jiku] (0,0) -- (0,2.8) node[above] {$x_3$};
  \draw[hojo] (0,0) rectangle (1.6,1.6);
  \draw[thick] (0,0) -- (1.6,0.48) -- (1.6,2.08) -- (0,1.6) -- cycle;
  \node[right] at (1.6,0.48) {$(\ell,\gamma\ell)$};
  \node[right] at (1.6,2.08) {$(\ell,(1+\gamma)\ell)$};
  \node[left] at (0,1.6) {$\ell$};
 \end{tikzpicture}
}
\kotae{\shoumonkai{1} $F_{32}=-F_{23}=\Omega_1/2$；$D_{23}=(\gamma-\Omega_1)/2$，$D_{32}=(\gamma+\Omega_1)/2$，他は0 \shoumonkai{2} $\Omega_1=\gamma$：$D_{32}=\gamma$，$\delta\vect{r}=(0,0,\gamma\,\delta x_2)$；$\Omega_1=-\gamma$：$D_{23}=\gamma$，$\delta\vect{r}=(0,\gamma\,\delta x_3,0)$ \shoumonkai{3} $x_2$ 軸に沿う辺が $x_3$ 側へ $\gamma$ 傾いた平行四辺形。$x_1$ 軸まわりに $x_2\to x_3$ の向きへ角度 $\gamma/2$}
\kenzan 逆向きの確認：$D_{32}=\gamma$ だけの $\tens{D}$ を分解すると $E_{23}=E_{32}=\gamma/2$，$F_{32}=\gamma/2$ で，\noteshiki{33} より $\Omega_1=2F_{32}=\gamma$。設定に戻る。また $\vect{\nabla}\times\vect{r}$ の第1成分
\[ \pdv{r_3}{x_2}-\pdv{r_2}{x_3}=\gamma \]
とも一致する。平行四辺形は底辺 $\ell$・高さ $\ell$ を保つので面積は変わらない。
\tsumazuki 回転角は $\Omega_1=\gamma$ ではなく $\gamma/2$。また $\Omega_1$ の符号で単純剪断のずれる向き（$x_3$ 方向か $x_2$ 方向か）が入れ替わる。
\end{kaisetsu}

\end{document}
